%!TEX root = ../main/main.tex

We now introduce the semantics for information extraction with errors. 
Let $S$ be a document spanner. 
Given a document $d$ and an error budget $k$, we consider edit scripts $E$ with $|E|\leq k$ and evaluate $S$ over the edited document $E(d)$. 
We write $\ValidEdits{k}{d}$ for the set of all edit scripts $E$ that are valid on $d$ and satisfy $|E|\leq k$.

\subsection{Approximate extraction semantics}

The semantics evaluates the spanner over an edited version of the input document, but the resulting mappings are transported backwards through the edit script to refer to valid spans within the original input document.

\begin{definition}
	Let $S$ be a document spanner, let $d \in \Sigma^*$, and let $k \in \mathbb{N}$. 
	The \emph{approximate semantics} of $S$ over $d$ with error budget $k$ is defined as
	\[
	\editsem{S}{k}{d} =
	\left\{
	(E^{-1}(\mu),E)
	\mid
	E \in \ValidEdits{k}{d} 
	\text{ and }
	\mu \in \llbracket S \rrbracket (E(d))
	\right\}.
	\]
\end{definition}

In this semantics, the script $E$ witnesses how the input document is modified to satisfy the spanner. While the mapping $\mu$ is computed directly on the resulting edited document $E(d)$, the returned mapping $E^{-1}(\mu)$ correctly refers to the coordinates of the original document $d$. 


\subsubsection{Examples}

To illustrate how the inverse transport adjusts the spans back to the original document, we present two scenarios: an insertion and a deletion occurring inside an extracted span.

\noindent
\textbf{Insertion inside an extracted span.}

Consider the original document $d=\mathtt{ac}$ and the edit script $E=\ins_{2,\mathtt{b}}$.
Applying the edit yields the document $E(d)=\mathtt{abc}$.
Suppose the spanner extracts the following mapping over the edited document:
\[
\mu =
[
x \mapsto \Span{1,4},
y \mapsto \Span{2,3}
].
\]
This corresponds to the extraction $x\{\mathtt{a}\,y\{\mathtt{b}\}\,\mathtt{c}\}$ over $\mathtt{abc}$.

Since the edit inserted $\mathtt{b}$ at position $2$, the inverse transport $E^{-1}$ acts as a deletion at position $2$. This reduces by one the upper limit of the spans that cross this position. 
Thus, transporting the mapping back to the original document yields:
\[
E^{-1}(\mu)
=
[
x \mapsto \Span{1,3},
y \mapsto \Span{2,2}
].
\]
In the coordinates of the original document $d=\mathtt{ac}$, the variable $x$ captures the entire document, while $y$ becomes an empty span between $\mathtt{a}$ and $\mathtt{c}$. We can visualize this effect as:
\[
\text{edited: } x\{\mathtt{a}\,y\{\mathtt{b}\}\,\mathtt{c}\}
\qquad
\leadsto
\qquad
\text{original: } x\{\mathtt{a}\,y\{\epsilon\}\,\mathtt{c}\}.
\]


\noindent
\textbf{Deletion inside an extracted span.}

Consider the original document $d=\mathtt{abbc}$ and the edit script $E=\del_2$.
Applying this deletion results in the edited document $E(d)=\mathtt{abc}$.
Assume the spanner extracts the mapping $\mu$ over $E(d)$ given by:
\[
\mu =
[
x \mapsto \Span{1,4},
y \mapsto \Span{3,4}
].
\]
This corresponds to the extraction $x\{\mathtt{a}\mathtt{b}\,y\{\mathtt{c}\}\}$ over $\mathtt{abc}$.

Because the original edit script deleted the symbol at position $2$ of $d$, the inverse transport $E^{-1}$ behaves as an insertion at position $2$. This effectively shifts the coordinates of the spans that occur after the deletion one position to the right, aligning them with the original text. 
Therefore, the inverse mapping is:
\[
E^{-1}(\mu)
=
[
x \mapsto \Span{1,5},
y \mapsto \Span{4,5}
].
\]
Over the original document $d=\mathtt{abbc}$, the variable $x$ correctly captures the entire string $\Span{1,5}$, and $y$ captures the final $\mathtt{c}$ at $\Span{4,5}$.

%!TEX root = ../main/main.tex

\subsection{Validity of the Approximate Semantics}

Since $\editsem{S}{k}{d}$ contains pairs of the form $(E^{-1}(\mu),E)$, it is not itself a span relation. Therefore, it does not directly define a document spanner. To obtain the actual extraction result, we project away the edit script.

\begin{definition}
	Let $S$ be a document spanner, let $d\in\Sigma^*$, and let $k\in\bbN$.
	The \emph{projected edit semantics} of $S$ over $d$ with error budget $k$ is defined as
	\[
	\projeditsem{S}{k}{d} =
	\left\{
		E^{-1}(\mu)
		\mid
		E\in\ValidEdits{k}{d}
		\text{ and }
		\mu\in\sem{S}{E(d)}
		\right\}.
	\]
\end{definition}

The projected edit semantics is the object that we study as a document spanner. For every input document $d$, it returns only mappings over spans of $d$.

\subsubsection{Document Spanner proof}

To ensure that $\projeditsem{S}{k}{d}$ acts as a valid extraction framework over the original document, we must prove that any mapping transported back through $E^{-1}$ only contains valid spans over $d$. We prove this by induction on the sequence of edit operations.

First, we show that a single inverse edit operation preserves span validity with respect to the pre-edited document.

\begin{lemma}\label{lemma:single_edit_validity}
	Let $d \in \Sigma^*$ be a document and $e \in \operatorname{Edits}$ an atomic edit operation. 
	If $s = \Span{i,j} \in \operatorname{span}(e(d))$, then $e^{-1}(s) \in \operatorname{span}(d)$.
\end{lemma}

\begin{proof}
	Let $n = |d|$. We analyze the three possible types of atomic edits applied to $d$ to verify that $e^{-1}(s)$ is a valid span over $d$. This means its transformed coordinates must fall strictly within the range $[1, n+1]$ and preserve the order $i \leq j$.
	
	\begin{itemize}
		\item \textbf{$e = \sub_{p,c}$:} 
		The length of the document does not change, so $|e(d)| = n$. 
		Thus, $s = \Span{i,j} \in \operatorname{span}(e(d))$ implies $1 \leq i \leq j \leq n+1$. 
		Since $e^{-1} = \sub_{p,c}$, the inverse transport leaves the span completely unchanged: $e^{-1}(s) = \Span{i,j}$. 
		Both bounds trivially remain valid over $d$.
		
		\vspace{0.3cm}
		
		\item \textbf{$e = \ins_{p,c}$:} 
		The edited document grows by one character, so $|e(d)| = n+1$. 
		Thus, $s = \Span{i,j} \in \operatorname{span}(e(d))$ implies $1 \leq i \leq j \leq n+2$. 
		A valid insertion on the original document $d$ occurs at position $1 \leq p \leq n+1$. The inverse transport is $e^{-1} = \del_p$. According to its definition, we analyze three sub-cases:
		\begin{itemize}
			\item \textit{Case $p < i \leq j$:} Both coordinates are strictly greater than $p$, so both decrease by one. The new upper bound is $j-1 \leq (n+2)-1 = n+1$. The new lower bound is $i-1 \geq p \geq 1$.
			\item \textit{Case $i \leq p < j$:} Only $j$ is strictly greater than $p$, so only $j$ decreases. The new upper bound becomes $j-1 \leq n+1$. The lower bound remains $i \geq 1$.
			\item \textit{Case $i \leq j \leq p$:} Neither coordinate changes. Since $p$ is bounded by the maximum valid insertion position on $d$ ($p \leq n+1$), we inherently have $j \leq p \leq n+1$. The lower bound remains $i \geq 1$.
		\end{itemize}
		In all three sub-cases, the new coordinates are bounded within $[1, n+1]$, wich mean $e^{-1}(s) \in \operatorname{span}(d)$.
		
		\vspace{0.3cm}
		\item \textbf{$e = \del_p$:} 
		The edited document shrinks by one character, so $|e(d)| = n-1$. 
		Thus, $s = \Span{i,j} \in \operatorname{span}(e(d))$ implies $1 \leq i \leq j \leq n$. 
		A valid deletion on the original document $d$ occurs at position $1 \leq p \leq n$. The inverse transport is $e^{-1} = \ins_{p,c}$. We again analyze three sub-cases:
		\begin{itemize}
			\item \textit{Case $p \leq i \leq j$:} Both coordinates are greater than or equal to $p$, so both increase by one. The new upper bound is $j+1 \leq n+1$. The new lower bound is $i+1 \geq p+1 \geq 2 > 1$.
			\item \textit{Case $i < p \leq j$:} Only $j$ increases. The new upper bound is $j+1 \leq n+1$. The lower bound remains $i \geq 1$.
			\item \textit{Case $i \leq j < p$:} Neither coordinate changes. Since $s$ was a valid span over $e(d)$, we know $j \leq n < n+1$. The lower bound remains $i \geq 1$.
		\end{itemize}
		In all three sub-cases, the new coordinates fall within $[1, n+1]$.
	\end{itemize}
	Since the limits are respected and the monotonic nature of the inverse operations preserves $i \leq j$, we conclude that $e^{-1}(s) \in \operatorname{span}(d)$.
\end{proof}

Now we extend this property to arbitrary edit scripts, proving that the returned mappings belong to the original document.

\begin{proposition}
	Let $d \in \Sigma^*$ be a document, $E = e_1e_2\cdots e_m \in \operatorname{Edits}^*$ an edit script, and $\mu$ a mapping over $E(d)$. 
	Then $E^{-1}(\mu)$ is a valid mapping over $d$.
\end{proposition}

\begin{proof}
	We proceed by induction on the length of the edit script, $m = |E|$.
	
	\textbf{Base case ($m=0$):} 
	If $E = \epsilon$, then $E(d) = d$. 
	The inverse transport is $\epsilon^{-1}(\mu) = \mu$. 
	Since $\mu$ is a mapping over $E(d)=d$, it is trivially a valid mapping over $d$.
	
	\textbf{Inductive step:} 
	Assume the proposition holds for all edit scripts of length $m-1$.
	We can decompose $E$ as $E = E' e_m$, where $E' = e_1 \cdots e_{m-1}$. 

	The edited document is $E(d) = e_m(E'(d))$, so $\mu \in \operatorname{e_m(E'(d))}$, the Lemma~\ref{lemma:single_edit_validity} guarantees that $e_m^{-1}(\mu)$ is a valid mapping over the intermediate document $E'(d)$.
	
	Since $e_m^{-1}(\mu)$ is a valid mapping over $E'(d)$ and $|E'| = m-1$, we apply the inductive hypothesis. 
	It follows that $(E')^{-1}(e_m^{-1}(\mu))$ is a valid mapping over $d$ and we have $E^{-1}(\mu) = (E')^{-1}(e_m(\mu))$. We conclude that $E^{-1}(\mu)$ is a valid mapping over $d$.
\end{proof}

Since every tuple in our semantics applies exactly this inverse transport, the approximate semantics correctly projects its extractions as valid mappings over the original document $d$.

\subsubsection{Regular spanner proof}

\begin{theorem}\label{thm:approx_regular_spanner}
Let $S$ be a regular document spanner and let $k\in\bbN$. Then $\projeditsem{S}{k}{\cdot}$ defines a regular document spanner.
\end{theorem}


\begin{proof}
	Since $S$ is a regular document spanner, there exists a variable-set automaton $\cA=(Q,\Sigma,V,q_0,F,\delta)$
	such that $\sem{\cA}{d}=\sem{S}{d}$ for every document $d\in\Sigma^*$.

	We construct a variable-set automaton $\mathcal{B}_k$ such that, for every document $d$,
	\[
	\sem{\mathcal{B}_k}{d}=\projeditsem{S}{k}{d}.
	\]
	The automaton $\mathcal{B}_k$ simulates $\cA$ while reading the original document. Its states store a state of $\cA$ and the number of edits used so far:
	\[
	Q' = Q \times \{0,\ldots,k\}.
	\]
	The initial state is $(q_0,0)$ and the final states are
	\[
	F' = F \times \{0,\ldots,k\}.
	\]
	
	The transition relation $\delta'$ is defined as follows. First, every non-editing transition of $\cA$ is copied without increasing the edit counter:
	\[
	(q,\alpha,q')\in\delta
	\quad\Longrightarrow\quad
	((q,r),\alpha,(q',r))\in\delta',
	\]
	for every $0\leq r\leq k$, where $\alpha$ is a letter, an $\varepsilon$-transition, or a variable marker.
	
	For every $0\leq r<k$, we also add the edit transitions. For every letter transition $(q,b,q')\in\delta$ and every $a\in\Sigma$ with $a\neq b$, we add
	\[
	((q,r),a,(q',r+1))\in\delta',
	\]
	which simulates substituting $a$ by $b$. For every $q\in Q$ and $a\in\Sigma$, we add
	\[
	((q,r),a,(q,r+1))\in\delta',
	\]
	which simulates deleting $a$ from the original document. Finally, for every letter transition $(q,b,q')\in\delta$, we add
	\[
	((q,r),\varepsilon,(q',r+1))\in\delta',
	\]
	which simulates inserting $b$ into the edited document.
	
	
	We use the following canonical convention for deletions. If a deletion and a variable marker occur at the same simulated boundary of $\cA$, then $\mathcal{B}_k$ performs the deletion before emitting the marker.
	
	\santiago{Esta parte del orden canonico hay que revisarlo bien. Es fundamental para que el automata funcione bien pero no se si es aplicable en la practica realmente}
	
	We now explain why the mappings produced by $\mathcal{B}_k$ are exactly the inverse transports defined above. Consider an accepting run $\rho$ of $\cA$ over an edited document $E(d)$, and let $\rho'$ be the corresponding run of $\mathcal{B}_k$ over the original document $d$. The run $\rho'$ follows the same simulated transitions as $\rho$, but uses edit transitions whenever the edit script $E$ inserts, deletes, or substitutes symbols.
	
	We use $j$ for the current input index of $\rho$ over the edited document $E(d)$, and $i$ for the current input index of $\rho'$ over the original document $d$. Thus, $j$ is a boundary in the edited document, while $i$ is a boundary in the original document.
	
	The variable-marker transitions of $\cA$ are copied into $\mathcal{B}_k$. Hence, if $\cA$ has a transition $(q,m,q')$, where $m$ is either $\ox$ or $\cx$, then $\mathcal{B}_k$ has the corresponding transition
	\[
	((q,r),m,(q',r))
	\]
	for every edit counter $r$.
	
	However, the index at which a marker is emitted is determined by the configuration where the transition is taken. If $\rho$ takes $(q,m,q')$ from configuration $(q,j)$, then $m$ is emitted at boundary $j$ of $E(d)$. The corresponding run $\rho'$ takes $((q,r),m,(q',r))$ from a configuration $((q,r),i)$, so $m$ is emitted at boundary $i$ of $d$.
	
	The edit transitions of $\mathcal{B}_k$ are designed so that, after simulating each atomic edit, the boundary $i$ is exactly the inverse transport of the boundary $j$ according to the definitions of $e^{-1}$ given above. We check this case by case.
	
	\begin{itemize}
		\item If the step is an ordinary match, then both runs consume one symbol. Thus the indices are updated as
		\[
		(i,j)\mapsto(i+1,j+1).
		\]
		No edit operation is added, so the inverse transport is the identity on this step.
		
		\item If the step is a substitution, then $\rho'$ consumes one symbol from $d$ and the simulated run $\rho$ consumes one symbol from $E(d)$. Again,
		\[
		(i,j)\mapsto(i+1,j+1).
		\]
		This agrees with the definition $(\sub_{\ell,c})^{-1}=\sub_{\ell,c}$, since substitutions do not change document length and therefore do not shift span boundaries.
		
		\item If the step is an insertion into the edited document, then $\rho$ consumes the inserted symbol, while $\rho'$ consumes no symbol from the original document. Hence
		\[
		(i,j)\mapsto(i,j+1).
		\]
		The edited boundary advances, but the original boundary stays fixed. This is exactly the inverse behavior of an insertion, since $(\ins_{\ell,c})^{-1}=\del_{\ell}$: a boundary that moves across an inserted symbol in $E(d)$ is transported back to the same boundary in $d$.
		
		\item If the step is a deletion from the original document, then $\rho'$ consumes one symbol from $d$, while the simulated run $\rho$ does not consume any symbol from $E(d)$. Hence
		\[
		(i,j)\mapsto(i+1,j).
		\]
		The original boundary advances, while the edited boundary stays fixed. This is exactly the inverse behavior of a deletion, since $(\del_{\ell})^{-1}=\ins_{\ell,c}$: transporting a boundary back through a deletion shifts it to the corresponding boundary in the original document. By the canonical convention above, if a deletion and a marker occur at the same simulated boundary, the deletion is performed before emitting the marker, so the marker is placed at the boundary prescribed by the inverse transport.
		
	\end{itemize}
	
	\noindent Now let $x\in V$, and suppose that the run $\rho$ assigns
	\[
	\mu(x)=\Span{j_1,j_2}
	\]
	over the edited document $E(d)$. This means that $\rho$ emits $\ox$ at boundary $j_1$ and $\cx$ at boundary $j_2$. Let the corresponding run $\rho'$ emit these markers at boundaries $i_1$ and $i_2$ of the original document. By the case analysis above, each atomic edit transition updates the pair of boundaries exactly as the corresponding inverse span transport. Therefore,
	\[
	\Span{i_1,i_2}
	=
	E^{-1}(\Span{j_1,j_2}).
	\]
	Hence, if $\mu'$ is the mapping produced by $\rho'$, then
	\[
	\mu'(x)=E^{-1}(\mu(x)).
	\]
	Since this holds for every variable $x\in\dom(\mu)$, we obtain
	\[
	\mu'=E^{-1}(\mu).
	\]
	Thus, $\mathcal{B}_k$ does not compute $\mu$ and then apply $E^{-1}$ afterwards; rather, its run emits the same variable markers as $\cA$ at the positions already transported back to the original document.
	
	We now finish the proof. The previous argument shows the following correspondence: accepting runs of $\mathcal{B}_k$ over $d$ are in one-to-one correspondence with pairs consisting of an edit script $E\in\ValidEdits{k}{d}$ and an accepting run of $\cA$ over $E(d)$. Moreover, if the run of $\cA$ produces a mapping $\mu$, then the corresponding run of $\mathcal{B}_k$ produces exactly the transported mapping $E^{-1}(\mu)$.
	
	Therefore, for every document $d$,
\[
\begin{aligned}
	\sem{\mathcal{B}_k}{d}
	&=
	\left\{
	E^{-1}(\mu)
	\;\middle|\;
	E\in\ValidEdits{k}{d}
	\text{ and }
	\mu\in\sem{\cA}{E(d)}
	\right\} \\
	&=
	\left\{
	E^{-1}(\mu)
	\;\middle|\;
	E\in\ValidEdits{k}{d}
	\text{ and }
	\mu\in\sem{S}{E(d)}
	\right\} \\
	&=
	\projeditsem{S}{k}{d}.
\end{aligned}
\]

	The second equality follows from $\sem{\cA}{d'}=\sem{S}{d'}$ for every document $d'$, and the last equality is exactly the definition of the projected edit semantics.
	
	Thus, $\projeditsem{S}{k}{\cdot}$ is definable by the variable-set automaton $\mathcal{B}_k$. Therefore, it is a regular document spanner.
	
\end{proof}
